% Section III — Static data do not determine the memory time
\section{Static data do not determine the memory time}\label{sec:identifiability}

Fix the harmonic vectors $\ee$ and $\jj\in\Hsp\setminus\{0\}$ and an operator class
$\mathfrak X\in\{\mathfrak P(\ee),\mathfrak E(Z)\}$. The question of this section is the size of the set
$\{\tmem(\Cop):\Cop\in\mathfrak X,\ \mathcal D(\Cop)=\mathcal D^\star\}$ for given static data
$\mathcal D^\star$. In Example~\ref{ex:four} and Theorem~\ref{thm:psd} rates are measured in units of a
common rate and times in its inverse.

\subsection{The positive-semidefinite class}\label{sec:psd}

\begin{example}[Same data, different memory; \status{proved}]\label{ex:four}
Let $n=4$, $\ee=(1,1,1,1)\tr/2$ and $\jj=(1,0,0,-1)\tr$. Let $\Cop_A=\tfrac43I-\tfrac13\mathbf 1\mathbf 1\tr$
and $\Cop_B=I-\tfrac12\mathsf A_4$, where $\mathbf 1=(1,1,1,1)\tr$ and $\mathsf A_4$ is the adjacency
matrix of the cycle $1\!-\!2\!-\!3\!-\!4\!-\!1$. Both operators have unit diagonal and kernel
$\mathrm{span}\{\ee\}$, and
\begin{equation}
K_A(p)=\frac{2}{p+4/3},\qquad K_B(p)=\frac1{p+1}+\frac1{p+2}.
\label{eq:four}
\end{equation}
Hence $K_A(0)=K_B(0)=3/2$, whereas $\tmem=3/4$ for $\Cop_A$ and $5/6$ for $\Cop_B$.
\end{example}

Both operators are weighted graph Laplacians on the four modes, so both lie in the single event cone
generated by the six vectors $\bm\delta_\mu-\bm\delta_\nu$. Static data therefore determine neither
the operator nor $\tmem$ even inside a fixed cone; Theorem~\ref{thm:cone} below shows that they bound
$\tmem$ there. In the positive-semidefinite class no bound exists.

\begin{theorem}[Unbounded memory at fixed static data; \status{proved}]\label{thm:psd}
Let $n=8$ and label the modes by $i\in\{0,\dots,7\}$. Let $\mathsf H$ be the symmetric orthogonal
Walsh matrix $\mathsf H_{ik}=8^{-1/2}(-1)^{\langle i,k\rangle}$, where $\langle i,k\rangle$ counts the
binary digits equal to one in both $i$ and $k$, and let $\mathbf h_k$ be its columns. For $0<\eta<1$ let
$\mathsf T(\eta)$ be the symmetric matrix whose only nonzero entries are
\begin{equation}
\begin{gathered}
\mathrm{diag}\,\mathsf T=(0,\eta^2,2,2-\eta^2,1,1,1,1),\\
\mathsf T_{12}=\mathsf T_{21}=\eta,\qquad \mathsf T_{47}=\mathsf T_{74}=-\eta ,
\end{gathered}
\label{eq:walsh}
\end{equation}
and put $\Cop_\eta=\mathsf H\,\mathsf T(\eta)\,\mathsf H\tr$, $\ee=\mathbf h_0$ and $\jj=\mathbf h_2$. Then:
\begin{enumerate}
\item[(i)] $\Cop_\eta\in\mathfrak P(\ee)$;
\item[(ii)] $\mathrm{diag}\,\Cop_\eta=(1,\dots,1)$;
\item[(iii)] $K_\eta(p)=(p+\eta^2)/[p^2+(2+\eta^2)p+\eta^2]$, hence $K_\eta(0)=1$ and
$\tmem(\Cop_\eta)=1+\eta^{-2}$;
\item[(iv)] the fixed-point-free involution $i\mapsto i\oplus3$ (bitwise exclusive or) commutes with
$\Cop_\eta$, leaves $\ee$ invariant and reverses $\jj$;
\item[(v)] for $|\omega|\ge\omega_{\min}>0$, $\big|K_\eta(i\omega)\,(i\omega+2)-1\big|\le\eta^2/(2\omega_{\min})$.
\end{enumerate}
Consequently, in $\mathfrak P(\ee)$ the static data
$(\mathbf r,\ker\Cop,K_0)=\big((1,\dots,1),\mathrm{span}\{\ee\},1\big)$ are compatible with every value
$\tmem\in(2,\infty)$.
\end{theorem}

\begin{proof}
(i) $\mathsf H$ is the character table of $(\mathbb Z_2)^3$ normalized to be orthogonal; it is symmetric.
$\mathsf T(\eta)$ is block diagonal on the index sets $\{0\}$, $\{1,2\}$, $\{3\}$, $\{4,7\}$, $\{5\}$,
$\{6\}$. The block on $\{1,2\}$ has determinant $\eta^2>0$ and positive trace; the block on $\{4,7\}$ has
determinant $1-\eta^2>0$; the remaining diagonal entries are $0$, $2-\eta^2>0$, $1$, $1$. Hence
$\mathsf T\succeq0$ with kernel $\mathrm{span}\{\bm\delta_0\}$, and $\Cop_\eta\succeq0$ with kernel
$\mathrm{span}\{\mathsf H\bm\delta_0\}=\mathrm{span}\{\ee\}$.
(ii) Since $\mathsf H_{ik}^2=1/8$ and $\mathrm{tr}\,\mathsf T=8$,
$(\Cop_\eta)_{ii}=1+2\eta(\mathsf H_{i1}\mathsf H_{i2}-\mathsf H_{i4}\mathsf H_{i7})$. Parity of
$\langle i,\cdot\rangle$ is additive under exclusive or, so
$\mathsf H_{i1}\mathsf H_{i2}=\tfrac18(-1)^{\langle i,3\rangle}=\mathsf H_{i4}\mathsf H_{i7}$, and the
bracket vanishes.
(iii) $(p+\Cop_\eta)^{-1}=\mathsf H(p+\mathsf T)^{-1}\mathsf H\tr$, so $K_\eta(p)$ is the $(2,2)$ entry of
the inverse of $p+\mathsf T$ restricted to $\{1,2\}$, which is $(p+\eta^2)/[(p+\eta^2)(p+2)-\eta^2]$. At
$p=0$ this is $1$; differentiating, $-K_\eta'(0)=1+\eta^{-2}$. Equivalently,
$\Cop_\eta(-\mathbf h_1/\eta+\mathbf h_2)=\jj$, so $\xx=-\mathbf h_1/\eta+\mathbf h_2$, $K_0=1$ and
$N=1+\eta^{-2}$.
(iv) Let $\mathsf P$ be the permutation $i\mapsto i\oplus3$. Then $\mathsf P\mathbf h_k=(-1)^{\langle3,k\rangle}\mathbf h_k$,
so $\mathsf P$ commutes with $\Cop_\eta$ if and only if $\mathsf T$ couples only indices of equal sign
$(-1)^{\langle 3,k\rangle}$. This holds: the signs are $-1$ for $k=1,2$ and $+1$ for $k=4,7$. The sign
is $+1$ for $k=0$ and $-1$ for $k=2$.
(v) $K_\eta(p)(p+2)-1=\eta^2/D(p)$ with $D(p)=p^2+(2+\eta^2)p+\eta^2$, and
$|D(i\omega)|\ge|\mathrm{Im}\,D(i\omega)|=(2+\eta^2)|\omega|\ge2|\omega|$.
For the consequence, take $\eta=(\tmem-1)^{-1/2}$.
\end{proof}

\begin{remark}[Mechanism]
$D$ has a root $-\eta^2/2+O(\eta^4)$ whose residue in $K_\eta$ is $\eta^2/4+O(\eta^4)$: a relaxation mode
of vanishing current weight and diverging lifetime carries about half of $K_0$ and all of the divergence
of $N$. Part~(v) shows that at frequencies $|\omega|\ge\omega_{\min}$ the response is within a relative
$\eta^2/(2\omega_{\min})$ of the single fast pole $1/(p+2)$. A band-limited measurement does not see
the divergence of $\tmem$; the slope at $p=0$ is not a property of the response in a finite band.
\end{remark}

\begin{remark}[Spectral reading and attribution]\label{rem:stieltjes}
By Eq.~\eqref{eq:stieltjes}, $|\jj|^2$, $K_0$ and $N$ are the moments of orders $0$, $-1$ and $-2$ of
the spectral measure of $\Cop|_{\Hsp}$ in the state $\jj$. The static data and the harmonic spectrum fix
the first two, and $\tmem$ requires the third. Without further structure no bound follows: this is the
elementary indeterminacy of a Stieltjes function from finitely many of its moments~\cite{Akhiezer2020,Schmudgen2017};
the same representation underlies classical bounds on effective transport coefficients of
composites~\cite{Golden1983}.
A spectral gap $\Gamma_{\min}$ gives $\tmem\le\Gamma_{\min}^{-1}$ (Lemma~\ref{lem:forms}). The content
of Theorem~\ref{thm:psd} is that prescribing the diagonal and the kernel does not supply such
structure; after a diagonal scaling, the positive-semidefinite matrices with prescribed diagonal form an
affine image of the elliptope~\cite{Saunderson2012}. The physical expectation is anticipated by the parity
structure of relaxons~\cite{Simoncelli2020} and by numerical failures of rate-only
descriptions~\cite{MalviyaRavichandran2023}; we have not found the identifiability statement itself for
phonon kinetics in the literature searched.
\end{remark}

\begin{remark}[The family is not microscopic]\label{rem:notcone}
No fixed finite event cone contains $\Cop_\eta$ for all small $\eta$. Let $\mathbf u=\mathbf h_1$; then
$\mathbf u\tr\Cop_\eta\jj=\eta$ and $\mathbf u\tr\Cop_\eta\mathbf u=\eta^2$. For
$\Cop=\sum_\gamma g_\gamma\bm\zeta_\gamma\bm\zeta_\gamma\tr$ with $g\ge0$,
\begin{equation}
|\mathbf u\tr\Cop\jj|\le\sum_\gamma g_\gamma|\mathbf u\tr\bm\zeta_\gamma|\,|\jj\tr\bm\zeta_\gamma|\le R\,\mathbf u\tr\Cop\mathbf u,
\label{eq:obstruction}
\end{equation}
with $R=\max\{|\jj\tr\bm\zeta_\gamma|/|\mathbf u\tr\bm\zeta_\gamma|:\mathbf u\tr\bm\zeta_\gamma\ne0\}$
(events with $\mathbf u\tr\bm\zeta_\gamma=0$ do not contribute). Membership for all small $\eta$ would
require $\eta\le R\eta^2$, which fails for $\eta<1/R$; if $R=0$ or no event has
$\mathbf u\tr\bm\zeta_\gamma\ne0$, the cross term vanishes, which contradicts $\eta>0$.
\end{remark}

\subsection{What static data do determine}\label{sec:lower}

\begin{proposition}[Lower bound from static data; \status{proved}]\label{prop:lower}
For every $\Cop\in\mathfrak P(\ee)$ and $\jj\in\Hsp\setminus\{0\}$,
\begin{equation}
\tmem\ \ge\ \frac{|\xx|}{|\jj|}\ \ge\ \tstat=\frac{K_0}{|\jj|^2},
\label{eq:lower}
\end{equation}
and $\tmem=\tstat$ if and only if $\jj$ is an eigenvector of $\Cop$.
\end{proposition}

\begin{proof}
By the Cauchy--Schwarz inequality $K_0=\jj\tr\xx\le|\jj|\,|\xx|$, so $|\xx|/K_0\ge1/|\jj|$ and
$\tmem=|\xx|\,(|\xx|/K_0)\ge|\xx|/|\jj|\ge K_0/|\jj|^2$. Equality throughout holds if and only if
$\xx\parallel\jj$, that is, $\Cop\jj\parallel\jj$.
\end{proof}

By Eq.~\eqref{eq:harmonic}, $\tstat=\lambda/\Lambda_i$ is the time of the elementary kinetic estimate
$\lambda=\Lambda_i\tau$; it is computable from the dc conductivity and the harmonic spectrum. In the
RTA it is the $c_\nu v_{\nu,i}^2$-weighted mean lifetime, and Proposition~\ref{prop:lower} reduces to
$\tmem^{\rm RTA}\ge\tstat^{\rm RTA}$. Equality holds in the grey Callaway model
(Proposition~\ref{prop:callaway}). Static data thus fix a lower end of the range of $\tmem$; whether that
end is approached in a given event geometry is a numerical question \PENDING{cross-reference: minima of
$\tmem/\tstat$ found in the Debye and AlN event geometries, Sec.~\ref{sec:aln}}.

The static data are not independent of each other. For an event operator, the dc response is bounded
below by the RTA response built from the same diagonal, up to the factor $1/\varsigma_{\max}$.

\begin{lemma}[Consistency of lifetimes and dc response; \status{proved}]\label{lem:rta}
Let $g$ be admissible, $\mathbf r=\mathrm{diag}\,\Cop(g)$ and $K_{\rm RTA}=\sum_\nu j_\nu^2/r_\nu$, and let
$\varsigma_{\max}$ be the largest number of distinct modes touched by an active event ($3$ for
$\mathsf p\to\mathsf a+\mathsf b$ with $\mathsf a\ne\mathsf b$, $2$ for $\mathsf p\to2\mathsf a$, up to $4$ for
four-phonon events). Then $K_0\ge K_{\rm RTA}/\varsigma_{\max}$. Equality holds for some
$\jj\in\Hsp\setminus\{0\}$ if and only if $\Cop(g)\yy=\varsigma_{\max}\,\mathrm{diag}(\mathbf r)\,\yy$ for some
$\yy\ne0$, and the currents attaining it are $\jj=\Cop(g)\yy$ with such $\yy$; this requires every active
event to touch exactly $\varsigma_{\max}$ modes. The constant is attained for $\varsigma_{\max}=2$, $3$ and $4$.
When some active event touches fewer than $\varsigma_{\max}$ modes, as for three-phonon events together with
four-phonon events on four modes, the inequality is strict, but the constant cannot be improved.
\end{lemma}

\begin{proof}
For $\jj\in\mathrm{range}\,\Cop$, $K_0=\max_{\yy}(2\jj\tr\yy-\yy\tr\Cop\yy)$, the variational principle of
the linearized Boltzmann equation~\cite{Ziman2001}. For a vector $\bm\zeta$ with $|\mathrm{supp}\,\bm\zeta|=s$,
the Cauchy--Schwarz inequality gives $(\bm\zeta\tr\yy)^2\le s\sum_\nu\zeta_\nu^2y_\nu^2$, that is
$\bm\zeta\bm\zeta\tr\preceq s\,\mathrm{diag}(\bm\zeta\circ\bm\zeta)$; summing with the rates,
$\Cop\preceq\varsigma_{\max}\,\mathrm{diag}\,\Cop$. With $\yy=\mathrm{diag}(\mathbf r)^{-1}\jj/\varsigma_{\max}$,
$K_0\ge2K_{\rm RTA}/\varsigma_{\max}-K_{\rm RTA}/\varsigma_{\max}$. Equality forces
$\yy\tr(\varsigma_{\max}\,\mathrm{diag}\,\Cop-\Cop)\yy=0$ for this $\yy$, hence
$\Cop\yy=\varsigma_{\max}\,\mathrm{diag}(\mathbf r)\,\yy=\jj$, because the matrix is positive semidefinite.
Conversely, if $\Cop\yy=\varsigma_{\max}\,\mathrm{diag}(\mathbf r)\,\yy$ with $\yy\ne0$, then
$\jj=\Cop\yy\in\Hsp$ is nonzero because all $r_\nu>0$, and
$K_0=\yy\tr\Cop\yy=\varsigma_{\max}\,\yy\tr\mathrm{diag}(\mathbf r)\,\yy=K_{\rm RTA}/\varsigma_{\max}$. Each active event contributes to
$\yy\tr(\varsigma_{\max}\,\mathrm{diag}\,\Cop-\Cop)\yy$ a nonnegative term, which vanishes if and only if the
numbers $\zeta_{\gamma\nu}y_\nu$, $\nu\in\mathrm{supp}\,\bm\zeta_\gamma$, are equal, and equal to zero when
$|\mathrm{supp}\,\bm\zeta_\gamma|<\varsigma_{\max}$. A zero of $\yy$ on one mode of an active event therefore
spreads to all its modes. The active events of an admissible network connect all modes, since a set of
modes not connected to the others would carry an energy invariant of its own; so one active event with
fewer than $\varsigma_{\max}$ modes forces $\yy=0$. Exact networks attain the ratio $1/\varsigma_{\max}$ for
$\varsigma_{\max}=2$, $3$ and $4$, and a family of mixed networks approaches it (Supplemental Material,
Sec.~S2).
\end{proof}

The lemma refers to $\mathrm{diag}\,\Cop$, in which the self-coupling of a repeated daughter is counted
in full. A self-energy linewidth that counts it with half weight is a different quantity, and the
inequality can fail for it (an exact counterexample is given in the Supplemental Material, Sec.~S2).

\subsection{Fixed event cone}\label{sec:cone}

\begin{definition}[Basic solutions and the constant $M$]\label{def:M}
For $I\subset\mathcal E$ with $|I|=n_{\Hsp}-1$, consider the linear system
\begin{equation}
\bm\zeta_\gamma\tr\yy=0\ \ (\gamma\in I),\qquad \yy\in\Hsp,\qquad \jj\tr\yy=1 .
\label{eq:basic}
\end{equation}
It has a unique solution $\yy_I$ if and only if $\{\bm\zeta_\gamma\}_{\gamma\in I}\cup\{\jj\}$ is a basis
of $\Hsp$; such $I$ are called bases, and $\mathfrak B$ denotes their set, which is nonempty because $Z$
spans $\Hsp$ and $\jj\ne0$. Put $M=M(Z,\jj)=\max_{I\in\mathfrak B}|\yy_I|$.
\end{definition}

$M$ does not depend on the rates. It depends on $Z$ only through the hyperplanes
$\mathrm{span}\{\bm\zeta_\gamma\}_{\gamma\in I}$, and $M(Z,s\jj)=M(Z,\jj)/|s|$, so that $M^2K_0$ and
$M^2|\jj|^2$ do not depend on the normalization of $\jj$.

\begin{theorem}[Cone bound; \status{proved}]\label{thm:cone}
Let $Z$ span $\Hsp$, $\jj\in\Hsp\setminus\{0\}$ and $g$ admissible. Fix an orthonormal basis of $\Hsp$;
for every $I\subset\mathcal E$ with $|I|=n_{\Hsp}-1$ define $\bm\psi_I\in\Hsp$ by $\bm\psi_I\tr\yy=\det[\bm\zeta_I,\yy]$ for all
$\yy\in\Hsp$, the determinant having the columns $\bm\zeta_\gamma$ ($\gamma\in I$, in a fixed order)
followed by $\yy$, and put $g_I=\prod_{\gamma\in I}g_\gamma$. Then
\begin{equation}
\frac{\xx}{K_0}=\sum_{I\in\mathfrak B}\varpi_I\,\yy_I,\qquad
\varpi_I=\frac{g_I\,(\bm\psi_I\tr\jj)^2}{\sum_{J\in\mathfrak B}g_J\,(\bm\psi_J\tr\jj)^2},
\label{eq:hull}
\end{equation}
a convex combination over the active bases ($g_I>0$). Consequently
\begin{equation}
|\Cop(g)^+\jj|\le M\,K_0(g),\qquad \tmem(g)\le M^2K_0(g).
\label{eq:conebound}
\end{equation}
\end{theorem}

\begin{proof}
Identify $\Hsp$ with $\mathbb R^{n_{\Hsp}}$ through the orthonormal basis and let
$A=\Cop(g)|_{\Hsp}=Z\,\mathrm{diag}(g)\,Z\tr$, where $Z$ also denotes the matrix with columns
$\bm\zeta_\gamma$; $A$ is positive definite. For $\yy,\mathbf z\in\Hsp$ the matrix determinant lemma gives
$\det(A+\yy\mathbf z\tr)=\det A+\mathbf z\tr\mathrm{adj}(A)\,\yy$. Writing
$A+\yy\mathbf z\tr=[Z,\yy]\,\mathrm{diag}(g,1)\,[Z,\mathbf z]\tr$ and expanding by the Cauchy--Binet formula
over the $n_{\Hsp}$-element subsets of columns, the subsets without the last column give $\det A$ and
those containing it give $\sum_{|I|=n_{\Hsp}-1}g_I\det[\bm\zeta_I,\yy]\det[\bm\zeta_I,\mathbf z]$. Hence
\begin{equation}
\mathrm{adj}(A)=\sum_{|I|=n_{\Hsp}-1}g_I\,\bm\psi_I\bm\psi_I\tr ,
\label{eq:adj}
\end{equation}
where dependent subsets give $\bm\psi_I=0$. With $A^{-1}=\mathrm{adj}(A)/\det A$,
$\xx=\sum_Ig_I\bm\psi_I(\bm\psi_I\tr\jj)/\det A$ and $K_0=\sum_Ig_I(\bm\psi_I\tr\jj)^2/\det A>0$. Now
$\bm\psi_I\tr\jj=\det[\bm\zeta_I,\jj]\ne0$ exactly when $I\in\mathfrak B$, and then
$\bm\psi_I/(\bm\psi_I\tr\jj)$ satisfies Eq.~\eqref{eq:basic}, because $\bm\psi_I\tr\bm\zeta_\gamma$ is a
determinant with a repeated column for $\gamma\in I$; it is therefore $\yy_I$. Dividing $\xx$ by $K_0$
gives Eq.~\eqref{eq:hull}. Finally $|\xx|/K_0\le\sum_I\varpi_I|\yy_I|\le M$ and
$\tmem=(|\xx|/K_0)^2K_0$.
\end{proof}

\begin{remark}[Relation to weighted least squares]\label{rem:bt}
By Lagrange's rule, $\xx/K_0$ minimizes $\sum_\gamma g_\gamma(\bm\zeta_\gamma\tr\yy)^2$ over $\yy\in\Hsp$
subject to $\jj\tr\yy=1$: it is the solution of a weighted least-squares problem with one exact
constraint. Ben-Tal and Teboulle proved, by the same Cauchy--Binet argument, that the solution of a
weighted least-squares problem lies in the convex hull of the solutions of its square subsystems for
every positive weight vector~\cite[Theorem~2.1 and Corollary~2.1]{BenTalTeboulle1990} (see also
Ref.~\cite{Berg1986}). Theorem~\ref{thm:cone} is that statement for the constrained problem, which is the
limit of infinite weight on the row $\jj\tr\yy=1$. Weight-independent bounds on scaled projections and
pseudoinverses are classical~\cite{Stewart1989,OLeary1990,Todd1990,HankeNeumann1993} and govern the
complexity of interior-point methods~\cite{VavasisYe1996}; approximating a closely related
weight-uniform condition measure is NP-hard in general~\cite{Tuncel1999}, so for large event sets $M$
must be bounded by exploiting structure. We claim no novelty for the inequality; its use to bound the
memory moment of a kinetic operator by its dc response is potentially novel within the literature
searched.
\end{remark}

The bound extends to rates that are not admissible but satisfy $\jj\in\mathrm{range}\,\Cop(g)$. Put
$g_\theta=g+\theta\mathbf 1$ ($\theta>0$), which is admissible, and $E=\Cop(\mathbf 1)$. Split
$\Hsp=\mathcal R\oplus\mathcal Q$ with $\mathcal R=\mathrm{range}\,\Cop(g)$ and
$\mathcal Q=\ker\Cop(g)\cap\Hsp$, and write $\xx_\theta=\Cop(g_\theta)^+\jj=\mathbf w_\theta+\mathbf u_\theta$
accordingly. The $\mathcal Q$ component of $\Cop(g_\theta)\xx_\theta=\jj$ gives
$\mathbf u_\theta=-(\Pi_{\mathcal Q}E\Pi_{\mathcal Q})^{-1}\Pi_{\mathcal Q}E\,\mathbf w_\theta$, and the
$\mathcal R$ component shows $\mathbf w_\theta\to\mathbf w=\Cop(g)^+\jj$ as $\theta\to0$; hence
$\mathbf u_\theta\to\mathbf u_\infty$ and $K_0(g_\theta)=\jj\tr\mathbf w_\theta\to K_0(g)$. Since
$\mathbf w\perp\mathbf u_\infty$, $|\mathbf w|\le|\mathbf w+\mathbf u_\infty|\le M K_0(g)$ by
Theorem~\ref{thm:cone} and continuity.

\begin{proposition}[Sharpness without lifetime constraints; \status{proved}]\label{prop:sharp}
For every $K^\star>0$,
\begin{equation}
\sup\{\tmem(g):\ g\ \text{admissible},\ K_0(g)=K^\star\}=M^2K^\star .
\label{eq:sharp}
\end{equation}
The supremum is approached by rates concentrated on a basis $I^\ast$ with $|\yy_{I^\ast}|=M$ and is in
general not attained; along the approaching sequence the rates on $I^\ast$ diverge, and so do the
relaxation rates of the modes touched by $I^\ast$.
\end{proposition}

\begin{proof}
For $0<\epsilon<1$ let $g^\epsilon_\gamma=1$ on $I^\ast$ and $\epsilon$ elsewhere; $g^\epsilon$ is
admissible. For $I\in\mathfrak B$, $g^\epsilon_I=\epsilon^{|I\setminus I^\ast|}$ with
$|I\setminus I^\ast|\ge1$ for $I\ne I^\ast$, so $\varpi_{I^\ast}\to1$ in Eq.~\eqref{eq:hull} and
$|\xx|/K_0\to M$ as $\epsilon\to0$. Since $\Cop(sg)=s\Cop(g)$, the ratio $|\xx|/K_0$ is invariant under
$g\mapsto sg$ while $K_0$ scales as $1/s$; rescaling $g^\epsilon$ to $K_0=K^\star$ gives
$\tmem=(|\xx|/K_0)^2K^\star\to M^2K^\star$, and Theorem~\ref{thm:cone} gives the reverse inequality.
Every $n_{\Hsp}$-element subset of events contains an event outside $I^\ast$, and some $I^\ast\cup\{\gamma\}$
with $\gamma\notin I^\ast$ is a basis of $\Hsp$, because the $\bm\zeta_\gamma$ ($\gamma\in I^\ast$) are
independent and $Z$ spans $\Hsp$; so $\det A=\Theta(\epsilon)$ while
$\sum_Ig^\epsilon_I(\bm\psi_I\tr\jj)^2\to(\bm\psi_{I^\ast}\tr\jj)^2>0$; thus $K_0(g^\epsilon)\sim c/\epsilon$
and the rescaling factor $K_0(g^\epsilon)/K^\star$ diverges.
\end{proof}

\begin{remark}[With lifetimes fixed]\label{rem:lifetimes}
Since $\mathcal F(\mathbf r^\star,K^\star)$ is contained in the set of Eq.~\eqref{eq:sharp},
$\sup_{\mathcal F}\tmem\le M^2K^\star$. Equality is possible: if $\jj=\bm\zeta_\beta$ is parallel to an
event vector, no subset containing $\beta$ is a basis, $M$ is attained by some $J^\ast\not\ni\beta$, and
every $g$ supported on the basis $J^\ast\cup\{\beta\}$ of $\Hsp$ gives $\xx/K_0=\yy_{J^\ast}$ exactly; with
$\mathbf r^\star=\mathrm{diag}\,\Cop(g)$ and $K^\star=K_0(g)$ the lifetime-constrained supremum is then
attained (an exact example with seven modes and ten events is given in the Supplemental Material,
Sec.~S3). Physical currents~\eqref{eq:current} vanish where $v_{\nu,i}=0$; they are odd under time
reversal, so their support consists of pairs $\{\nu,\bar\nu\}$ of modes of equal energy, and they are never
parallel to a three-phonon event vector, whose support has three modes or two modes of different energies.
For them it is \status{open} whether $\sup_{\mathcal F}\tmem<M^2K^\star$. For rates tied by a symmetry
group, the proof of Proposition~\ref{prop:sharp} transfers when some maximizing $I^\ast$ is a union of
orbits; otherwise the supremum over symmetric rates is not known.
\end{remark}

For three-phonon events, $M$ has a reading in terms of sub-networks.

\begin{proposition}[Sub-network form of $M$; \status{proved}]\label{prop:subnetwork}
Let $\bm\zeta_\gamma=\mathsf W^{-1}\mathbf s_\gamma$ as in Eq.~\eqref{eq:events}, and for $I\in\mathfrak B$
let $\mathsf S_I$ be the matrix with columns $\mathbf s_\gamma$, $\gamma\in I$. Then
$\ker\mathsf S_I\tr=\mathrm{span}\{\bm\varepsilon,\bm\phi_I\}$ for some $\bm\phi_I\notin\mathrm{span}\{\bm\varepsilon\}$, and
\begin{equation}
\yy_I=\frac{\Pi_{\Hsp}\mathsf W\bm\phi_I}{\jj\tr\mathsf W\bm\phi_I},\qquad
|\yy_I|=\frac{|\Pi_{\Hsp}\mathsf W\bm\phi_I|}{|\jj\tr\mathsf W\bm\phi_I|}.
\label{eq:subnetwork}
\end{equation}
\end{proposition}

\begin{proof}
Since $\mathsf W$ is invertible, the $\bm\zeta_\gamma$ ($\gamma\in I$) are independent if and only if
$\mathsf S_I$ has full column rank $n-2$; then $\dim\ker\mathsf S_I\tr=2$, and $\bm\varepsilon\in\ker\mathsf S_I\tr$
by energy conservation. The conditions $\bm\zeta_\gamma\tr\yy=\mathbf s_\gamma\tr\mathsf W^{-1}\yy=0$
($\gamma\in I$) say $\mathsf W^{-1}\yy\in\ker\mathsf S_I\tr$, i.e.,
$\yy=a\,\mathcal V^{1/2}\ee+b\,\mathsf W\bm\phi_I$. The condition $\yy\in\Hsp$ gives
$\yy=b\,\Pi_{\Hsp}\mathsf W\bm\phi_I$, and $\jj\tr\yy=1$ with $\jj\perp\ee$ gives
$b=1/(\jj\tr\mathsf W\bm\phi_I)$; the denominator is nonzero because $I\in\mathfrak B$.
\end{proof}

The vector $\bm\phi_I$ is a conserved quantity of the events in $I$ other than energy,
$\mathbf s_\gamma\tr\bm\phi_I=0$ for $\gamma\in I$. $M$ is large when some $(n-2)$-event sub-network
has such a hidden invariant with a small overlap with the current. By Eq.~\eqref{eq:hull},
$|\xx|/K_0\le\sum_I\varpi_I|\yy_I|$, so a large ratio $\tmem/K_0$ requires convex weight on bases with
large $|\yy_I|$: for $0<\epsilon\le1/2$, $\tmem\ge(1-\epsilon)^2M^2K_0$ puts at least half of the weight on
bases with $|\yy_I|\ge(1-2\epsilon)M$. The weight resides in the products $g_I(\bm\psi_I\tr\jj)^2$, not
necessarily in individual rates. Proposition~\ref{prop:sharp} shows that concentrating the rates on one
such basis attains the bound, which turns the invariant of the sub-network into a slow, weakly
current-coupled relaxation mode, the mechanism of Theorem~\ref{thm:psd}.

\begin{corollary}[Width of the feasible range; \status{proved}]\label{cor:gap}
If $\mathcal F(\mathbf r^\star,K^\star)$ is nonempty, then
$\sup_{\mathcal F}\tmem\,/\,\inf_{\mathcal F}\tmem\le M^2|\jj|^2$.
\end{corollary}

\begin{proof}
Combine $\tmem\le M^2K^\star$ (Theorem~\ref{thm:cone}) with $\tmem\ge K^\star/|\jj|^2$
(Proposition~\ref{prop:lower}).
\end{proof}

The structure of the feasible set is convex in part. With the extended value
$K_0(g)=\sup_{\yy}\{2\jj\tr\yy-\yy\tr\Cop(g)\yy\}$, equal to $\jj\tr\Cop(g)^+\jj$ if
$\jj\in\mathrm{range}\,\Cop(g)$ and $+\infty$ otherwise, $K_0$ is convex, lower semicontinuous and
nonincreasing in $g$, with $\partial K_0/\partial g_\gamma=-(\bm\zeta_\gamma\tr\xx)^2$ at admissible $g$.
By the generalized Schur complement,
\begin{multline}
\{g\ge0:\ \mathrm{diag}\,\Cop(g)=\mathbf r^\star,\ K_0(g)\le K^\star\}\\
=\Big\{g\ge0:\ \mathrm{diag}\,\Cop(g)=\mathbf r^\star,\
\begin{pmatrix}\Cop(g)&\jj\\ \jj\tr&K^\star\end{pmatrix}\succeq0\Big\},
\label{eq:lmi}
\end{multline}
so the infimum of $K_0$ at fixed lifetimes is a semidefinite program (\status{proved}). The set contains
nonadmissible rates; the infimum over admissible rates coincides with it whenever the lifetime
constraint admits an admissible point. No convexity of $\tmem$ is claimed.

\subsection{What bounds the memory: amplitude information}\label{sec:prior}

\begin{proposition}[Per-event prior; \status{proved}]\label{prop:prior}
Let $g^{\rm ref}$ be admissible and $F>0$. Every $g$ with $g_\gamma\ge g^{\rm ref}_\gamma/F$ for all
$\gamma\in\mathcal E$ is admissible and satisfies
\begin{equation}
\tmem(g)\le\frac{F}{\Gamma_{\min}\big(\Cop(g^{\rm ref})\big)},
\label{eq:prior}
\end{equation}
where $\Gamma_{\min}$ is the smallest eigenvalue on $\Hsp$. No lifetime or dc data enter, and the bound
holds with or without symmetry constraints on the rates. Conversely, for $\tau>0$, $\tmem(g)>\tau$ requires
$\min\{g_\gamma/g^{\rm ref}_\gamma:g^{\rm ref}_\gamma>0\}<[\tau\,\Gamma_{\min}(\Cop(g^{\rm ref}))]^{-1}$.
\end{proposition}

\begin{proof}
$\Cop(g)-F^{-1}\Cop(g^{\rm ref})=\sum_\gamma(g_\gamma-g^{\rm ref}_\gamma/F)\bm\zeta_\gamma\bm\zeta_\gamma\tr\succeq0$,
so $\Gamma_{\min}(\Cop(g))\ge\Gamma_{\min}(\Cop(g^{\rm ref}))/F>0$, which implies admissibility, and
Lemma~\ref{lem:forms} gives $\tmem(g)\le1/\Gamma_{\min}(\Cop(g))$. For the converse, if $g_\gamma=0$ for
some $\gamma$ with $g^{\rm ref}_\gamma>0$, the minimum is $0$ and there is nothing to prove; otherwise
$g\ge g^{\rm ref}/F$ with $F=\max\{g^{\rm ref}_\gamma/g_\gamma:g^{\rm ref}_\gamma>0\}$, and the bound gives
the converse.
\end{proof}

The information that confines $\tmem$ is therefore microscopic: a statement about every event rate
relative to a reference, which is what an anharmonic amplitude model provides and what static data do
not. The quantitative consequences for AlN are reported in Sec.~\ref{sec:aln}.

\subsection{Infrared criterion}\label{sec:ir}

We now let the mesh be refined. Consider uniform $\Gamma$-centered meshes $\mathbf k=\mathsf B\mathbf n/N_1$
($\mathbf n\in\mathbb Z^d$, $\mathsf B$ the matrix of primitive reciprocal-lattice vectors, so that
$|\det\mathsf B|\,V_{\rm cell}=(2\pi)^d$) with $N_1$ points per reciprocal-lattice direction in spatial
dimension $d$ ($N_{\mathbf k}=N_1^d$), the $\Gamma$ point excluded. With the normalization~\eqref{eq:current}, mode sums $\sum_\nu F(\mathbf k_\nu,s_\nu)/\mathcal V$
are Riemann sums for $(2\pi)^{-d}\sum_s\int_{\rm BZ}F\,d^d\mathbf k$. Write $j_\nu=\mathcal V^{-1/2}J(\mathbf k_\nu,s_\nu)$
with $J=w\varepsilon v_i$, and define the RTA moments
\begin{equation}
\mathcal M_m=\sum_s\int_{\rm BZ}J(\mathbf k,s)^2\,r(\mathbf k,s)^{-m}\,\frac{d^d\mathbf k}{(2\pi)^d},\qquad m=0,1,2.
\label{eq:Mm}
\end{equation}
On the meshes, $|\jj|^2$, $K_{\rm RTA}$ and $\tmem^{\rm RTA}$ are $\mathcal M_0(N_1)$, $\mathcal M_1(N_1)$
and $\mathcal M_2(N_1)/\mathcal M_1(N_1)$, with the mesh moments
$\mathcal M_m(N_1)=\mathcal V^{-1}\sum_\nu J(\mathbf k_\nu,s_\nu)^2r_\nu^{-m}$; their continuum counterparts are
$\mathcal M_0$, $\mathcal M_1$ and $\mathcal M_2/\mathcal M_1$.
On each acoustic branch, with $k=|\mathbf k|$ and $\hat{\mathbf k}=\mathbf k/k$, the hypotheses are:
\begin{enumerate}
\item[(H1)] $r(\mathbf k)=r_0(\hat{\mathbf k})\,k^\alpha[1+o(1)]$ as $k\to0$, uniformly in $\hat{\mathbf k}$, with
$0<r_-\le r_0(\hat{\mathbf k})\le r_+<\infty$;
\item[(H2)] $J(\mathbf k)\to J_0(\hat{\mathbf k})$ as $k\to0$, uniformly in $\hat{\mathbf k}$; $J$ is bounded on the
Brillouin zone; and $0<A_m=\sum_{\rm ac}\int_{S^{d-1}}J_0^2\,r_0^{-m}\,d\Omega<\infty$ for the orders $m$ considered;
\item[(H3)] away from $\mathbf k=0$, and on optical branches, $r$ is bounded below by a positive constant.
\end{enumerate}
For a branch with linear dispersion and sound velocity $v_s$, $w\varepsilon\to k_BT$ and
$J_0^2=(k_BTv_s)^2\hat k_i^2$, which vanishes on the plane $\hat k_i=0$; (H2) requires only that the
angular integral be positive. For the event operator one more hypothesis is needed:
\begin{enumerate}
\item[(H4)] for the admissible operators $\Cop_{N_1}$ on the meshes and $\xx=\Cop_{N_1}^+\jj$, there are
$c>0$ and $k_0>0$, independent of $N_1$, such that $|x_\nu|\ge c\,|j_\nu|/r_\nu$ for every mode with
$k_\nu<k_0$.
\end{enumerate}
Limits on the meshes require, in addition, some regularity:
\begin{enumerate}
\item[(H5)] for the orders $m$ considered, $\Phi_m=\sum_{\rm ac}J_0^2r_0^{-m}$ is bounded on $S^{d-1}$ and
continuous outside a closed set of directions of zero measure, and on every branch $J^2r^{-m}$ is Riemann
integrable on the part $k\ge\delta$ of the Brillouin zone for every $\delta>0$.
\end{enumerate}
Hypothesis (H5) admits the discontinuities of frequency-sorted acoustic branches at the acoustic axes,
where they touch, and where they cross; some regularity of this kind is necessary (Supplemental Material,
Sec.~S4).

\begin{theorem}[Infrared criterion]\label{thm:ir}
Let the meshes, the moments $\mathcal M_m$ and hypotheses (H1)--(H5) be as above.
\begin{enumerate}
\item[(i)] \emph{RTA moments} (\status{proved}). Assume (H1)--(H3). Then $\mathcal M_m<\infty$ if and
only if $m\alpha<d$; for $m\alpha=d$ the moment restricted to $k>k_{\min}$ diverges like $\log(1/k_{\min})$.
In the RTA the conductivity is therefore finite if and only if $\alpha<d$, and the memory moment if and
only if $2\alpha<d$; for $\alpha\ge d$, where $\mathcal M_1$ and $\mathcal M_2$ both diverge, the ratio of
the moments restricted to $k>k_{\min}$ diverges like $k_{\min}^{-\alpha}$ for $\alpha>d$ and like
$k_{\min}^{-d}/\log(1/k_{\min})$ for $\alpha=d$. If, in addition, (H5) holds, then
$\mathcal M_m(N_1)\to\mathcal M_m$ for $m\alpha<d$,
$\mathcal M_m(N_1)=(2\pi)^{-d}A_m\log N_1+o(\log N_1)$ for $m\alpha=d$, and
$\mathcal M_m(N_1)=N_1^{m\alpha-d}S_m/V_{\rm cell}\,[1+o(1)]$ for $m\alpha>d$, where
$S_m=\sum_{\mathbf n\ne0}\Phi_m(\mathsf B\mathbf n/|\mathsf B\mathbf n|)\,|\mathsf B\mathbf n|^{-m\alpha}\in(0,\infty)$.
Hence on the meshes $\tmem^{\rm RTA}$ converges to $\mathcal M_2/\mathcal M_1$ for $2\alpha<d$ and grows
like $\log N_1$ for $2\alpha=d$, like $N_1^{2\alpha-d}$ for $d/2<\alpha<d$, and like $N_1^d/\log N_1$ for
$\alpha=d$.
\item[(ii)] \emph{Event operator} (\status{derived under assumptions}, namely (H4) and the uniform form of
(H1)). Suppose that (H2) holds, that $\Phi_2$ is continuous outside a closed set of directions of zero
measure, that $\mathbf r=\mathrm{diag}\,\Cop_{N_1}$ satisfies (H1) uniformly in $N_1$, that is, with $r_0$
independent of $N_1$ and the $o(1)$ uniform in $\hat{\mathbf k}$ and in $N_1$, and that (H4) holds. If
$2\alpha\ge d$, then $N=|\Cop_{N_1}^+\jj|^2\to\infty$ and $\tmem\ge N^{1/2}/|\jj|\to\infty$ as
$N_1\to\infty$, with no condition on $K_0$; for $2\alpha>d$, $\tmem\gtrsim N_1^{\alpha-d/2}$, and for
$2\alpha=d$, $\tmem\gtrsim(\log N_1)^{1/2}$. These are lower bounds; two-sided mesh exponents for the
event operator do not follow.
\item[(iii)] \emph{Cone constant} (\status{derived under assumptions}, namely (H4) and the uniform form
of (H1)). Under the hypotheses of (ii) and $2\alpha\ge d$, $M_{N_1}^2K_{0,N_1}\to\infty$, where $M_{N_1}$ is the constant of
Theorem~\ref{thm:cone} for the event family on the mesh; if, in addition, $K_{0,N_1}$ stays bounded,
$M_{N_1}\to\infty$.
\end{enumerate}
\end{theorem}

\begin{proof}
(i) By (H3) and the boundedness of $J$, the integrand of Eq.~\eqref{eq:Mm} is bounded outside a ball
$k<k_0$. Inside it, (H1)--(H2) give
$J^2r^{-m}=[J_0(\hat{\mathbf k})^2r_0(\hat{\mathbf k})^{-m}+\epsilon(\mathbf k)]\,k^{-m\alpha}$ with
$\sup_{\hat{\mathbf k}}|\epsilon|\to0$ as $k\to0$; the error is additive because $J_0$ may vanish on a
plane. In polar coordinates $d^d\mathbf k=k^{d-1}dk\,d\Omega$ the ball contributes
$\int_{k_{\min}}^{k_0}k^{d-1-m\alpha}[A+\bar\epsilon(k)]\,dk$, where $A=A_m>0$ and $\bar\epsilon(k)\to0$.
Choosing $k_0$ with $|\bar\epsilon|\le A/2$ bounds this between $A/2$ and $3A/2$ times
$\int_{k_{\min}}^{k_0}k^{d-1-m\alpha}dk$, which stays bounded as $k_{\min}\to0$ if and only if $m\alpha<d$,
equals $\log(k_0/k_{\min})$ for $m\alpha=d$ and is of order $k_{\min}^{d-m\alpha}$ for $m\alpha>d$; the
statement for $\alpha\ge d$ follows. For the mesh laws, the optical branches and the part $k\ge k_0$ of
the acoustic branches contribute a bounded amount to $\mathcal M_m(N_1)$, by (H3) and the boundedness of
$J$. On the acoustic
branches the part $k_\nu<k_0$ equals
$N_1^{m\alpha-d}V_{\rm cell}^{-1}\sum_{0<|\mathsf B\mathbf n|<k_0N_1}[\varphi(\mathbf n)+\epsilon_{\rm ac}(\mathsf B\mathbf n/N_1)\,|\mathsf B\mathbf n|^{-m\alpha}]$,
where $\varphi(\mathbf n)=\Phi_m(\mathsf B\mathbf n/|\mathsf B\mathbf n|)\,|\mathsf B\mathbf n|^{-m\alpha}$ and
$\epsilon_{\rm ac}$ is the sum of $\epsilon$ over the acoustic branches. Lemma~S1 of the Supplemental
Material (Sec.~S4) evaluates such sums. On each dyadic shell $2^l\le|\mathsf B\mathbf n|<2^{l+1}$ the sum
of $\varphi$ is $2^{l(d-m\alpha)}$ times a Riemann sum of the homogeneous extension of $\varphi$ over the
shell $1\le|\mathsf B\mathbf u|<2$, which converges to a positive limit by (H5); the $\epsilon_{\rm ac}$
terms are bounded shell by shell. Summing over the shells, with $|\det\mathsf B|\,V_{\rm cell}=(2\pi)^d$, gives
the stated laws for $m\alpha\ge d$. For $m\alpha<d$ the part $k<\delta$ of $\mathcal M_m(N_1)$ and of
$\mathcal M_m$ is $O(\delta^{d-m\alpha})$ uniformly in $N_1$, and the remaining sums are Riemann sums that
converge by (H5), so that $\mathcal M_m(N_1)\to\mathcal M_m$. Taking the ratio $\mathcal M_2(N_1)/\mathcal M_1(N_1)$ gives the four
cases.
(ii) By (H4), $N\ge\sum_{k_\nu<k_0}x_\nu^2\ge c^2\sum_{k_\nu<k_0}j_\nu^2r_\nu^{-2}$. Only the upper side
of the uniform (H1) is used: $r_\nu\le r_0(\hat{\mathbf k}_\nu)k_\nu^\alpha[1+o(1)]$ with an $o(1)$
independent of $N_1$, so the last sum is bounded below by a mesh sum of the form treated in part~(i),
with angular function $\Phi_2$ and an error independent of $N_1$; by Lemma~S1 it grows like
$N_1^{2\alpha-d}$ for $2\alpha>d$ and like $\log N_1$ for $2\alpha=d$. Since $|\jj|^2$ is at most the
number of branches times $\sup J^2/V_{\rm cell}$, Proposition~\ref{prop:lower} gives
$\tmem\ge N^{1/2}/|\jj|\to\infty$ at the stated rates.
(iii) Theorem~\ref{thm:cone} gives $\tmem\le M_{N_1}^2K_{0,N_1}$ on every mesh.
\end{proof}

\begin{example}[Finite dc response, infinite slope; \status{proved}]\label{ex:ir}
The Stieltjes function
$K(p)=\int_0^1q^2(q^2+p)^{-1}dq=1-\sqrt p\,\arctan(1/\sqrt p)$ is the RTA response of a one-dimensional
branch with $J^2\propto q^2$ and $r\propto q^2$. It has $K(0)=1$ and $K(p)=1-\tfrac\pi2\sqrt p+O(p)$: the
small-$p$ expansion is not analytic and no Cattaneo time exists. With a cutoff $q>1/N_1$,
$K_0=1-1/N_1$ and $\tmem=N_1$.
\end{example}

Hypothesis (H4) is a property of a given family of operators, not of a class: it states that the
long-wavelength part of the solution $\xx$ is RTA-like, and it holds with $c=1$ in the RTA. In three
dimensions a finite conductivity requires $\alpha<3$, whereas a finite memory moment requires
$\alpha<3/2$. The low-frequency three-phonon rates of acoustic phonons follow power laws whose exponents
depend on the branch and on the crystal symmetry~\cite{Herring1954}; whether the operators of a given
material satisfy (H1)--(H5), and with which $\alpha$, is a separate question \PENDING{AlN: infrared law
of the acoustic rates and status of (H1)--(H5) for the AlN operators, Sec.~\ref{sec:aln}}. The mesh laws
are approached slowly: in the lattice test of Supplemental Sec.~S4, $\mathcal M_1(N_1)$ converges like
$N_1^{-(d-\alpha)}$ for $d/2<\alpha<d$, and at $\alpha=d$ the relative correction decays like
$1/\log N_1$, so that exponents fitted on coarse meshes are biased. When the
memory moment diverges, a Cattaneo time exists only for a regularized problem: a finite sample, whose
boundaries scatter long-wavelength phonons, or a finite frequency band. This motivates the
film-regularized memory time of Sec.~\ref{sec:aln}.
